\begin{explanationbox}
\textbf{\textcolor{CExplanation}{一句话理解}}
把圆沿一个方向拉长，就得到椭圆；圆的“半径垂直切线”变成椭圆的
“径向斜率 $\times$ 切线斜率 $=-b^2/a^2$”。
\textbf{不再等于 $-1$，是因为横、纵两个方向被拉伸的倍数不同。}

\medskip
\textbf{\textcolor{CExplanation}{一、先从最熟悉的圆出发}}
圆上取点 $Q$，连接圆心与 $Q$，圆的切线与半径垂直。
中学里两条斜率存在的直线互相垂直，等价于 $k_1k_2=-1$。
现在把单位圆 $u^2+v^2=1$ 按
\[
        x=au,\qquad y=bv\qquad(a>b>0)
\]
拉伸，得到椭圆 $x^2/a^2+y^2/b^2=1$。
\textbf{任何一条非竖直直线的斜率} $\Delta v/\Delta u$，拉伸后都变成
\[
  \frac{\Delta y}{\Delta x}
  =\frac{b\,\Delta v}{a\,\Delta u}
  =\frac ba\frac{\Delta v}{\Delta u}.
\]
因此，圆上“半径”和“切线”的斜率\emph{各自}乘上 $b/a$，
两者乘积就由 $-1$ 变成
\[
   \left(\frac ba\right)^2(-1)=-\frac{b^2}{a^2}.
\]
这就是本结论最直观的来历：\textbf{两个方向都改变，所以比例要平方。}

\medskip
\textbf{\textcolor{CExplanation}{二、为什么斜率相乘后会是常数？}}
若 $P(x_0,y_0)$ 位于椭圆第一象限，$OP$ 向右上方走，
所以 $k_{OP}=y_0/x_0>0$；椭圆在那里向右下方弯，
切线 $k_l<0$。移动切点时，两条线都会转动：
一条越来越陡，另一条就按相反的比例变平。
\[
 k_{OP}=\frac{y_0}{x_0},\qquad
 k_l=-\frac{b^2}{a^2}\frac{x_0}{y_0}.
\]
乘起来，$x_0$ 和 $y_0$ 约掉，所以与点的位置无关。
这里的“径向”只是指从\textbf{椭圆中心 $O$} 连向 $P$，
\textbf{不是指从焦点连向 $P$}。

\medskip
\textbf{\textcolor{CExplanation}{三、看得见的图形记忆}}
想象 $a=4,b=3$：水平拉得比竖直长，椭圆显得扁。
圆的斜率乘积是 $-1$；拉伸后两条斜率分别按 $3/4$ 缩小，故
\[
 k_{OP}k_l=-\frac34\cdot\frac34=-\frac9{16}.
\]
椭圆越扁，$b/a$ 越小，乘积越接近 $0$（但始终小于 $0$）。
又因 $e^2=1-b^2/a^2$，因此 $-b^2/a^2=e^2-1$。
\textbf{速记：圆是负一，椭圆乘一个“短轴/长轴”的平方。}

\medskip
\textbf{\textcolor{CExplanation}{四、从已知切线方程一眼看出}}
已知椭圆在 $P$ 处的切线
\[
  \frac{x_0x}{a^2}+\frac{y_0y}{b^2}=1,
\]
其 $x$ 系数为 $x_0/a^2$，$y$ 系数为 $y_0/b^2$。
直线 $Ax+By=C$ 的斜率是 $-A/B$，所以
\[
    k_l=-\frac{x_0/a^2}{y_0/b^2}
       =-\frac{b^2x_0}{a^2y_0}.
\]
而 $k_{OP}=y_0/x_0$，两式相乘立即结束。
\textbf{这比考试现场重新求导更快。}

\medskip
\textbf{\textcolor{CExplanation}{五、与几何中的法向量联系}}
对于直线 $Ax+By=C$，向量 $(A,B)$ 垂直于直线，称为法向量。
因此切线的法向量可取
\[
 \boldsymbol n=\left(\frac{x_0}{a^2},\frac{y_0}{b^2}\right).
\]
而 $\overrightarrow{OP}=(x_0,y_0)$。
圆上 $a=b$ 时，$\boldsymbol n$ 与 $\overrightarrow{OP}$ 平行，
所以 $OP$ 垂直切线；椭圆上 $a\ne b$ 时，
横纵分量被\emph{不同倍数} 缩放，通常不再平行。
\textbf{这说明椭圆的“法线”一般不通过中心，不能把椭圆当圆。}

\medskip
\textbf{\textcolor{CExplanation}{六、考试识别信号}}
题目一旦同时出现“中心 $O$、椭圆上一点 $P$、$OP$ 的斜率、
$P$ 处切线斜率、离心率”，先想
\[
 \boxed{k_{OP}k_l=e^2-1.}
\]
若仅给切点坐标并要求整条切线方程，优先写切点式；
若给切线斜率、要反求离心率，优先用乘积式；
若点在坐标轴上，改用切线方程，不写斜率乘积。
\end{explanationbox}
